\documentclass[11pt]{article}
\usepackage{series}
\title{Measurement Effects as Finite Survivor--Basin Kernels\\
\large Projective Collapse as the Singular Limit of Admissible Selection in Modal Triplet Theory}
\author{Peter Nero}
\date{April 2026}
\begin{document}
\maketitle

\begin{abstract}
Ideal quantum measurement is commonly represented by projection-valued measures, sharp position eigenstates, and Dirac delta kernels. In the preceding papers of this sequence, Dirac delta functions were reinterpreted as singular shadows of admissible projection, coherent Green functions replaced point sources by finite kernels, and gauge fixing was re-read as admissible section selection. The present paper applies the same projection-kernel program to measurement.

The central claim is narrow: ideal projective measurement is the zero-width limit of a finite admissible measurement effect. Mathematically, the rank-one distributional projector \(|x\rangle\langle x|\), whose kernel is \(\delta(y-x)\delta(z-x)\), is replaced by a finite coherent effect
\[
E_x^{\mathrm{coh}}(y,z)=K_{\mathrm{det}}(y,x)K_{\mathrm{det}}^\ast(z,x),
\]
or, more generally, by a positive operator-valued measure generated by a detector response kernel. In the sharp limit, this finite effect converges weakly to the usual projection-valued measurement.

In MTT language, measurement is not an additional primitive collapse rule. It is a localized disturbance followed by damping, projection, and stabilization into an admissible survivor basin. The measured outcome is the stabilized basin label; the Dirac delta is the singular notation obtained when the basin width, detector response, and residual disturbance floor are idealized away. We prove the finite-resolution POVM construction, its normalization, its sharp-limit convergence, the corresponding instrument/state update, and a discrete two-basin model. We then connect this to fixed-point stabilization, Ornstein--Uhlenbeck floors, and the broader delta diagnostic principle.
\end{abstract}

\section{Purpose and claim discipline}

This paper is the measurement sequel to the delta/projection sequence. The guiding principle is:
\[
\boxed{\text{ideal projective collapse is the zero-width limit of finite admissible selection}.}
\]

The paper does not claim to replace the operational formalism of quantum mechanics. It proves a narrower mathematical fact: finite detector-response kernels define positive measurement effects whose sharp limit reproduces ordinary projective measurement. The MTT interpretation is then that these finite effects are the downstream representation of survivor-basin selection after disturbance and re-coherence.

\subsection*{Non-claims}

We do not claim:
\begin{enumerate}
\item that all measurement devices are modeled by one universal kernel;
\item that the Born rule is rederived from first principles in this paper;
\item that standard projective measurements are invalid;
\item that detector physics is ignored;
\item that finite effects alone solve every interpretive issue in quantum mechanics.
\end{enumerate}

The narrower claim is:
\[
\boxed{
|x\rangle\langle x|
\text{ is an idealized zero-width measurement effect; the admissible object is finite-width.}
}
\]

\section{Ideal projective measurement and the delta kernel}

On \(L^2(X)\), the formal position eigenstate \(|x\rangle\) satisfies
\[
\langle y|x\rangle=\delta(y-x).
\]
The corresponding ideal rank-one projector is
\[
P_x=|x\rangle\langle x|.
\]
Its formal integral kernel is
\[
P_x(y,z)=\delta(y-x)\delta(z-x).
\]
This object is distributional rather than a bounded operator on \(L^2(X)\). More precisely, \(|x\rangle\) is not an element of \(L^2(X)\), and \(|x\rangle\langle x|\) is not a bounded rank-one operator on the Hilbert space. It is a useful rigged-Hilbert-space or distributional shorthand. Operationally, real detectors never implement \(P_x\) at zero width. They implement finite-resolution effects centered near \(x\).

This distinction is one of the main rigor advantages of the present formulation: for every \(\epsilon>0\), the finite detector effects below are honest bounded positive operators, while the sharp object is recovered only as a weak or distributional limit.

\section{Detector response kernels}

Let \(X\) be a smooth configuration space with measure \(d\mu\). Let \(r_\epsilon(x,y)\) be a detector response amplitude centered at \(x\), with resolution scale \(\epsilon>0\). Assume:
\begin{enumerate}
\item \(r_\epsilon(x,\cdot)\in L^2(X)\) for each \(x\);
\item \(r_\epsilon\) is measurable in both variables;
\item the family is normalized as below.
\end{enumerate}

Define
\[
(M_x^\epsilon\psi)(y)=r_\epsilon(x,y)\psi(y),
\qquad
E_x^\epsilon=(M_x^\epsilon)^\ast M_x^\epsilon.
\]
Then
\[
(E_x^\epsilon\psi)(y)=|r_\epsilon(x,y)|^2\psi(y).
\]

\begin{definition}[Finite detector POVM]
A family \(\{E_x^\epsilon\}_{x\in X}\) is a finite detector POVM if
\[
E_x^\epsilon\ge 0,
\qquad
\int_X E_x^\epsilon\,d\mu(x)=I
\]
in the weak operator sense.
\end{definition}

A sufficient normalization condition is
\[
\int_X |r_\epsilon(x,y)|^2\,d\mu(x)=1
\qquad
\text{for almost every }y\in X.
\]

\begin{theorem}[Finite detector effects define a POVM]
Assume \(r_\epsilon\) satisfies
\[
\int_X |r_\epsilon(x,y)|^2\,d\mu(x)=1
\]
for almost every \(y\). Then \(E_x^\epsilon=(M_x^\epsilon)^\ast M_x^\epsilon\) defines a positive operator-valued measure:
\[
E_x^\epsilon\ge 0,
\qquad
\int_XE_x^\epsilon\,d\mu(x)=I
\]
weakly on \(L^2(X)\).
\end{theorem}

\begin{proof}
For any \(\psi\in L^2(X)\),
\[
\langle \psi,E_x^\epsilon\psi\rangle
=
\int_X |r_\epsilon(x,y)|^2|\psi(y)|^2\,d\mu(y)\ge 0.
\]
Thus \(E_x^\epsilon\ge 0\). Moreover,
\[
\left\langle\psi,\left(\int_XE_x^\epsilon\,d\mu(x)\right)\psi\right\rangle
=
\int_X\int_X |r_\epsilon(x,y)|^2|\psi(y)|^2\,d\mu(x)d\mu(y).
\]
Using the normalization condition gives
\[
\left\langle\psi,\left(\int_XE_x^\epsilon\,d\mu(x)\right)\psi\right\rangle
=
\|\psi\|^2.
\]
By polarization the identity holds weakly. Hence the family is a POVM.
\end{proof}

\section{Sharp-limit convergence}

Let
\[
w_\epsilon(x,y):=|r_\epsilon(x,y)|^2.
\]
Assume \(w_\epsilon(x,y)\) is an approximate identity in the \(x\)-variable:
\[
w_\epsilon(x,y)\to\delta(x-y)
\]
distributionally as \(\epsilon\downarrow0\). The probability density for outcome \(x\) is
\[
p_\epsilon(x|\psi)
=
\langle \psi,E_x^\epsilon\psi\rangle
=
\int_X w_\epsilon(x,y)|\psi(y)|^2\,d\mu(y).
\]

\begin{theorem}[Sharp detector limit]
If \(w_\epsilon(x,y)\) is an approximate identity and \(|\psi|^2\) is continuous, then
\[
p_\epsilon(x|\psi)\to |\psi(x)|^2
\]
pointwise at continuity points of \(|\psi|^2\). In distributional form,
\[
p_\epsilon(\cdot|\psi)\to |\psi(\cdot)|^2.
\]
\end{theorem}

\begin{proof}
The expression
\[
p_\epsilon(x|\psi)
=
\int_X w_\epsilon(x,y)|\psi(y)|^2\,d\mu(y)
\]
is the smoothing of \(|\psi|^2\) by an approximate identity. Standard approximate-identity convergence gives pointwise convergence at continuity points and distributional convergence in general.
\end{proof}

Thus the Born density appears as the zero-width limit of finite detector response:
\[
\boxed{
|\psi(x)|^2
=
\lim_{\epsilon\downarrow0}
\langle\psi,E_x^\epsilon\psi\rangle.
}
\]

\section{Post-measurement update and instruments}

A POVM determines outcome probabilities, but a measurement model also requires a state-update rule. The standard object for this is an instrument. For the finite detector response above, define the completely positive operation
\[
\mathcal I_x^\epsilon(\rho)
=
M_x^\epsilon\rho (M_x^\epsilon)^\ast .
\]
The probability density of outcome \(x\) is
\[
p_\epsilon(x|\rho)
=
\operatorname{Tr}\mathcal I_x^\epsilon(\rho)
=
\operatorname{Tr}(\rho E_x^\epsilon),
\]
and the normalized post-measurement state is
\[
\rho\mapsto
\rho_x^\epsilon
=
\frac{\mathcal I_x^\epsilon(\rho)}
{\operatorname{Tr}\mathcal I_x^\epsilon(\rho)}.
\]

For a pure state \(\rho=|\psi\rangle\langle\psi|\), this reduces to
\[
\psi\mapsto
\psi_x^\epsilon
=
\frac{M_x^\epsilon\psi}{\|M_x^\epsilon\psi\|}.
\]
Explicitly,
\[
\psi_x^\epsilon(y)
=
\frac{r_\epsilon(x,y)\psi(y)}
{\left(\int_X |r_\epsilon(x,z)|^2|\psi(z)|^2\,d\mu(z)\right)^{1/2}}.
\]
This is not a collapse to a point. It is a finite filtering operation.

\begin{remark}[Formal square-root notation]
The expression \(\delta^{1/2}\) is not a standard distribution. It only indicates that the effect density \(w_\epsilon=|r_\epsilon|^2\) approaches \(\delta\). The finite-\(\epsilon\) measurement operators and instruments are the mathematically meaningful objects.
\end{remark}

\section{Gaussian model}

On \(X=\mathbb R^d\), a standard detector model is
\[
w_\epsilon(x,y)
=
(2\pi\epsilon^2)^{-d/2}
\exp\left(-\frac{|x-y|^2}{2\epsilon^2}\right).
\]
Then
\[
p_\epsilon(x|\psi)
=
(2\pi\epsilon^2)^{-d/2}
\int_{\mathbb R^d}
\exp\left(-\frac{|x-y|^2}{2\epsilon^2}\right)
|\psi(y)|^2\,dy.
\]
Thus the observed probability density is the Gaussian-smoothed Born density.

The post-measurement state is
\[
\psi_x^\epsilon(y)
=
\frac{
(2\pi\epsilon^2)^{-d/4}
\exp\left(-\frac{|x-y|^2}{4\epsilon^2}\right)\psi(y)
}{
p_\epsilon(x|\psi)^{1/2}
}.
\]
This is a finite-width localized survivor state, not an exact point eigenstate.


\section{Discrete outcomes and survivor-basin effects}

Position measurement is the cleanest place to display the delta kernel, but many physical measurements have discrete outcomes: spin up/down, detector clicks, threshold events, channel labels, or pointer states. The finite-effect formulation extends directly.

Let \(\mathcal H\) be a Hilbert space and let \(\{E_i\}_{i=1}^N\) be positive operators satisfying
\[
E_i\ge 0,
\qquad
\sum_{i=1}^N E_i=I.
\]
Then \(\{E_i\}\) is a discrete POVM. In the present interpretation, each \(E_i\) represents the finite downstream effect associated with a survivor basin \(B_i\).

If each \(E_i\) is a sharp projection \(P_i\) and
\[
P_iP_j=\delta_{ij}P_i,
\qquad
\sum_iP_i=I,
\]
then the measurement is projective. The finite-survivor-basin view reverses the usual hierarchy: the \(E_i\) are native, while the \(P_i\) arise when basins are idealized as perfectly separated.

\subsection{Two-basin qubit model}

A minimal two-basin model is a noisy or finite-resolution measurement of \(\sigma_z\). Define
\[
E_+
=
\frac12(I+\eta\sigma_z),
\qquad
E_-
=
\frac12(I-\eta\sigma_z),
\qquad
0\le \eta\le 1.
\]
The eigenvalues of \(E_\pm\) are nonnegative for \(0\le \eta\le1\), and
\[
E_+ + E_- = I.
\]
Thus \(\{E_+,E_-\}\) is a valid two-outcome POVM.

For a state
\[
\rho=\frac12(I+\vec r\cdot\vec\sigma),
\]
the probabilities are
\[
p_\pm
=
\operatorname{Tr}(\rho E_\pm)
=
\frac12(1\pm \eta r_z).
\]
When \(\eta=1\), this is the sharp projective measurement of \(\sigma_z\). When \(0<\eta<1\), the two outcome basins are distinguishable but not infinitely sharp. When \(\eta=0\), the device carries no information about the basin label.

Thus \(\eta\) is a basin-separation or selection-sharpness parameter:
\[
\boxed{
\eta=1 \text{ gives ideal projective selection; }
0\le\eta<1 \text{ gives finite survivor-basin selection.}
}
\]

A corresponding Lüders-type finite update may be written
\[
\mathcal I_\pm(\rho)
=
E_\pm^{1/2}\rho E_\pm^{1/2}.
\]
This gives a concrete discrete analogue of the continuous detector kernel: the finite measurement is an honest POVM/instrument pair, and the projective rule is recovered as the singular sharp-basin limit.


\section{MTT interpretation: survivor-basin kernels}

In MTT, measurement is an ordinary episode of the evolve--project cycle:
\[
T_t=\Pi_{\mathrm{coh}}\circ\Phi_t.
\]
The measuring apparatus introduces a localized disturbance. The subsequent dynamics damps incoherent components and projects the system back into the coherent/admissible sector. The outcome is the stabilized basin into which the system falls.

Thus the finite effect \(E_x^\epsilon\) is interpreted as the downstream measurement representation of a survivor basin:
\[
\boxed{
E_x^\epsilon
=
\text{finite downstream effect associated with a stabilized survivor basin}.
}
\]

The projective measurement \(P_x=|x\rangle\langle x|\) is the singular limit:
\[
P_x
=
\lim_{\epsilon\downarrow0}E_x^\epsilon
\]
only in the weak/distributional sense appropriate to idealized sharp measurements.

\section{OU floors and irreducible width}

The fixed-point disturbance analysis suggests that finite width is not merely instrumental. Persistent disturbances balanced by damping produce an Ornstein--Uhlenbeck floor:
\[
\sigma^2=\frac{\delta}{2\gamma},
\]
where \(\delta\) is disturbance strength and \(\gamma\) is the effective damping margin.

Thus a physically admissible measurement kernel should not generally take \(\epsilon\to0\). Instead, its effective width should satisfy a lower bound of the form
\[
\epsilon_{\mathrm{eff}}^2
\gtrsim
\frac{\delta}{2\gamma}.
\]
This gives an MTT-native refinement:
\[
\boxed{
\text{collapse width cannot be taken below the disturbance--damping floor of the regime}.
}
\]

\section{Relation to the delta diagnostic}

The present paper realizes the delta diagnostic in measurement theory. The formal kernel of exact projection onto a position eigenstate contains
\[
\delta(y-x)\delta(z-x).
\]
The finite replacement is a detector/survivor-basin effect. More generally,
\[
E_x^\epsilon = (M_x^\epsilon)^\ast M_x^\epsilon.
\]

Thus:
\[
\boxed{
\text{measurement delta}
=
\text{zero-width shadow of finite detector/survivor-basin effect}.
}
\]

\section{Comparison with standard POVM theory}

The finite effects used here are standard POVMs. The new claim is interpretive and structural: in an MTT setting, such POVMs are not merely practical approximations to ideal projectors. They are closer to the physically native object. The ideal projection-valued measure is the limiting fiction.

This reverses the usual hierarchy:
\[
\text{standard view: PVM fundamental, POVM practical};
\]
\[
\text{MTT view: finite POVM native, PVM singular idealization}.
\]

\section{Scope of proof}

\begin{center}
\begin{tabular}{p{0.28\linewidth}p{0.62\linewidth}}
\hline
\textbf{Status} & \textbf{Claim}\\
\hline
Proved & Finite detector kernels define POVMs under normalization.\\
Proved & Approximate-identity detector effects converge to sharp projective probabilities.\\
Proved & Gaussian finite-resolution measurement gives smoothed Born density and finite post-measurement filtering.\\
Proved & Discrete finite-outcome POVMs model imperfect survivor-basin separation, with projective measurement recovered in the sharp limit.\\
Interpretive & MTT identifies these finite effects with survivor-basin selection after disturbance and re-coherence.\\
Not proved here & Full derivation of the Born rule from basin volumes in the underlying modal theory.\\
Not proved here & A universal microscopic detector model for all measurement regimes.\\
\hline
\end{tabular}
\end{center}

\section{Consequences}

\begin{enumerate}
\item Ideal collapse is a singular limit, not the native measurement operation.
\item Finite measurement width is expected, not merely experimental imperfection.
\item Detector response, environmental disturbance, and basin geometry all contribute to the effective kernel.
\item Exact position eigenstates and perfectly separated discrete pointer states are downstream idealizations.
\item The measurement problem is softened: selection is stabilization into a basin, not a primitive discontinuous law.
\end{enumerate}

\section{Conclusion}

This paper extends the delta/projection program to measurement. The formal projector \(|x\rangle\langle x|\) is a distributional idealization whose kernel contains Dirac deltas. A physically meaningful measurement is represented by finite effects generated by detector response kernels. These effects form a POVM, reproduce projective measurement in the sharp limit, and yield finite post-measurement states.

In MTT language, the finite measurement effect is the downstream encoding of survivor-basin selection after disturbance and re-coherence. The delta appears only when basin width, detector response, and disturbance floors are idealized away.

The central result is therefore:
\[
\boxed{
\text{projective measurement}
=
\text{singular limit of finite admissible survivor-basin selection}.
}
\]

\end{document}
